\section{相关研究与方法定位}
\subsection{可复用技能与技能网络}
早期自主智能体通过可执行程序积累技能，Voyager 展示了技能库在持续探索中的复用作用\citep{P11}。
ReAct、Reflexion、Toolformer、ToolLLM、ToolkenGPT、LATM 与 HuggingGPT 分别从行动推理、反馈、工具学习及任务分解等角度改进能力调用\citep{P12,P13,P14,P15,P16,P17,P18}。
这些工作支持将技能视为可复用的行为资源，但工具可被调用与工具在具体任务上有效仍是不同问题。
现代开放技能生态进一步使说明文档与资源包本身成为评价对象。
SkillNet、SkillCT 与 SkillLearnBench 是本文的近邻工作；本文不把多维度、能力边界或分层评价视为首创\citep{P01,P43,P44}。

SkillsBench 强调固定环境、有无技能对照及独立任务验收\citep{P02}。
ASSAY 通过干预估计任务条件下的技能效应，提示一个全局好坏分数不能替代条件化贡献\citep{P03}。
SkillRouter 研究大规模候选路由，DSR 研究相关性与集合多样性的联合选择\citep{P04,P05}。
其中重排中的 quality 是相关性权重，不能被解释为技能安全性或内在可信度。
本文的前五检索对照采用词法表示，ECP使用开发产物证据与显式前提；均不是这些训练模型的完整复现。

\subsection{评审、证据与风险控制}
G-Eval、MT-Bench 与 Prometheus 2 提供了结构化模型评审方法\citep{P24,P25,P26}，
JudgeBench 与位置偏差研究则表明，评审稳定性和表面一致仍需要客观参照\citep{P27,P28}。
ALCE、FActScore 和 RARR 将事实拆分、来源支持与修订联系起来\citep{P39,P40,P41}。
本文借鉴其证据思想，但将“文件与行号真实”“语义支持成立”“产物正确”作为三个不同测量层次。
引用可定位只解决第一层，不继承另外两层的准确性。

开放技能中的风险还包括恶意指令和正常功能受损。
生态实证、AgentDojo、InjecAgent 与 ToolEmu 提供了攻击与正常任务联合评测的依据\citep{P09,P10,P29,P30,P31}。
本文不执行恶意技能代码，也不由静态线索推断恶意。
共形风险控制与风险控制预测集合为错误放行约束提供了统计工具\citep{P36,P37}；
但当前数据的来源集中、独立校准标签缺失，不满足可直接宣称分布无关保证的条件。

\subsection{更新与系统优化}
DSPy、TextGrad、AFlow、ADAS 和 GEPA 表明，结构化反馈、候选搜索与程序优化可以改进系统配置\citep{P19,P20,P21,P22,P23}。
Skill++、COBRA-Skills 与 SkillHEX 将元数据演化、预算分配和失败修补用于技能生态\citep{P06,P07,P08}。
这类方法的效果仍应由独立任务和回归评测决定。
DPP 提供集合选择的标准工具，LinUCB 为预算内探索提供条件化决策视角\citep{P38,P42}。
SWE-bench、AgentBench、ToolSandbox 与 $\tau$-bench 则强调不同任务环境下的验收与行为轨迹\citep{P32,P33,P34,P35}。
本文的增量定位为：在固定包版本上连接可反驳声明、独立产物判别、条件证据与完整兼容路径，
同时披露该连接尚未获得的外部效度。
